% Options for packages loaded elsewhere
% Options for packages loaded elsewhere
\PassOptionsToPackage{unicode}{hyperref}
\PassOptionsToPackage{hyphens}{url}
%
\documentclass[
  ignorenonframetext,
  aspectratio=169,
  russian,
  aspectratio=169]{beamer}
\newif\ifbibliography
\usepackage{pgfpages}
\setbeamertemplate{caption}[numbered]
\setbeamertemplate{caption label separator}{: }
\setbeamercolor{caption name}{fg=normal text.fg}
\beamertemplatenavigationsymbolshorizontal
% remove section numbering
\setbeamertemplate{part page}{
  \centering
  \begin{beamercolorbox}[sep=16pt,center]{part title}
    \usebeamerfont{part title}\insertpart\par
  \end{beamercolorbox}
}
\setbeamertemplate{section page}{
  \centering
  \begin{beamercolorbox}[sep=12pt,center]{section title}
    \usebeamerfont{section title}\insertsection\par
  \end{beamercolorbox}
}
\setbeamertemplate{subsection page}{
  \centering
  \begin{beamercolorbox}[sep=8pt,center]{subsection title}
    \usebeamerfont{subsection title}\insertsubsection\par
  \end{beamercolorbox}
}
% Prevent slide breaks in the middle of a paragraph
\widowpenalties 1 10000
\raggedbottom
\AtBeginPart{
  \frame{\partpage}
}
\AtBeginSection{
  \ifbibliography
  \else
    \frame{\sectionpage}
  \fi
}
\AtBeginSubsection{
  \frame{\subsectionpage}
}
\usepackage{iftex}
\ifPDFTeX
  \usepackage[T1]{fontenc}
  \usepackage[utf8]{inputenc}
  \usepackage{textcomp} % provide euro and other symbols
\else % if luatex or xetex
  \usepackage{unicode-math} % this also loads fontspec
  \defaultfontfeatures{Scale=MatchLowercase}
  \defaultfontfeatures[\rmfamily]{Ligatures=TeX,Scale=1}
\fi
\usepackage{lmodern}

\usefonttheme{serif} % use mainfont rather than sansfont for slide text
\ifPDFTeX\else
  % xetex/luatex font selection
  \setmainfont[]{DejaVu Serif}
  \setsansfont[]{DejaVu Sans}
  \setmonofont[]{DejaVu Sans Mono}
\fi
% Use upquote if available, for straight quotes in verbatim environments
\IfFileExists{upquote.sty}{\usepackage{upquote}}{}
\IfFileExists{microtype.sty}{% use microtype if available
  \usepackage[]{microtype}
  \UseMicrotypeSet[protrusion]{basicmath} % disable protrusion for tt fonts
}{}


\usepackage{longtable,booktabs,array}
\newcounter{none} % for unnumbered tables
\usepackage{calc} % for calculating minipage widths
\usepackage{caption}
% Make caption package work with longtable
\makeatletter
\def\fnum@table{\tablename~\thetable}
\makeatother
\usepackage{graphicx}
\makeatletter
\newsavebox\pandoc@box
\newcommand*\pandocbounded[1]{% scales image to fit in text height/width
  \sbox\pandoc@box{#1}%
  \Gscale@div\@tempa{\textheight}{\dimexpr\ht\pandoc@box+\dp\pandoc@box\relax}%
  \Gscale@div\@tempb{\linewidth}{\wd\pandoc@box}%
  \ifdim\@tempb\p@<\@tempa\p@\let\@tempa\@tempb\fi% select the smaller of both
  \ifdim\@tempa\p@<\p@\scalebox{\@tempa}{\usebox\pandoc@box}%
  \else\usebox{\pandoc@box}%
  \fi%
}
% Set default figure placement to htbp
\def\fps@figure{htbp}
\makeatother



\ifLuaTeX
\usepackage[bidi=basic,shorthands=off,provide=*]{babel}
\else
\usepackage[bidi=default,shorthands=off,provide=*]{babel}
\fi
\ifPDFTeX
\else
\babelfont{rm}[]{DejaVu Serif}
\fi
\ifLuaTeX
  \usepackage{selnolig} % disable illegal ligatures
\fi


\setlength{\emergencystretch}{3em} % prevent overfull lines

\providecommand{\tightlist}{%
  \setlength{\itemsep}{0pt}\setlength{\parskip}{0pt}}



 

\usepackage[]{csquotes}

\IfFileExists{plex-otf.sty}{
  %% Full TeXlive
  % \usepackage[%
  %   % math,
  %   RM={Scale=0.94},SS={Scale=0.94},SScon={Scale=0.94},TT={Scale=MatchLowercase,FakeStretch=0.9},DefaultFeatures={Ligatures=Common}
  % ]{plex-otf}
}{
  %% TinyTeX
  \usepackage{libertine}
}

%%% Load theme
% https://deic.uab.cat/~iblanes/beamer_gallery/
\IfFileExists{beamerthemegotham.sty}{
  %% Full TeXlive
  \usetheme{gotham}
  \gothamset{
    numbering=totalpagenumber,
    parttocframe default=off,
    sectiontocframe default=off,
    subsectiontocframe default=off,
  }
}{
  %% TinyTeX
  \usetheme{Madrid}
}
\makeatletter
\@ifpackageloaded{caption}{}{\usepackage{caption}}
\AtBeginDocument{%
\ifdefined\contentsname
  \renewcommand*\contentsname{Содержание}
\else
  \newcommand\contentsname{Содержание}
\fi
\ifdefined\listfigurename
  \renewcommand*\listfigurename{Список иллюстраций}
\else
  \newcommand\listfigurename{Список иллюстраций}
\fi
\ifdefined\listtablename
  \renewcommand*\listtablename{Список таблиц}
\else
  \newcommand\listtablename{Список таблиц}
\fi
\ifdefined\figurename
  \renewcommand*\figurename{Рисунок}
\else
  \newcommand\figurename{Рисунок}
\fi
\ifdefined\tablename
  \renewcommand*\tablename{Таблица}
\else
  \newcommand\tablename{Таблица}
\fi
}
\@ifpackageloaded{float}{}{\usepackage{float}}
\floatstyle{ruled}
\@ifundefined{c@chapter}{\newfloat{codelisting}{h}{lop}}{\newfloat{codelisting}{h}{lop}[chapter]}
\floatname{codelisting}{Список}
\newcommand*\listoflistings{\listof{codelisting}{Листинги}}
\makeatother
\makeatletter
\makeatother
\makeatletter
\@ifpackageloaded{caption}{}{\usepackage{caption}}
\@ifpackageloaded{subcaption}{}{\usepackage{subcaption}}
\makeatother

\usepackage{bookmark}
\IfFileExists{xurl.sty}{\usepackage{xurl}}{} % add URL line breaks if available
\urlstyle{same}
\hypersetup{
  pdftitle={Выполнение лабораторной работы №5},
  pdfauthor={Карпухин Клим},
  pdflang={ru-RU},
  hidelinks,
  pdfcreator={LaTeX via pandoc}}


\title{\texorpdfstring{Выполнение лабораторной работы
№5}{Выполнение лабораторной работы №5}}
\subtitle{\texorpdfstring{Настройка рабочей
среды}{Настройка рабочей среды}}
\author{Карпухин Клим}
\date{2026-03-14}

\begin{document}
\frame{\titlepage}


\section{Содержание}\label{ux441ux43eux434ux435ux440ux436ux430ux43dux438ux435}

\begin{frame}{Содержание}
\begin{enumerate}[<+->]
\tightlist
\item
  Информация о докладчике
\item
  Вводная часть и актуальность
\item
  Объект и предмет исследования
\item
  Научная новизна и практическая значимость
\item
  Цель, гипотеза и задачи
\item
  Материалы, методы и инструменты
\item
  Ход работы (этапы, скриншоты)
\item
  Результаты и анализ
\item
  Выводы
\end{enumerate}
\end{frame}

\section{Информация}\label{ux438ux43dux444ux43eux440ux43cux430ux446ux438ux44f}

\begin{frame}{Докладчик}
\protect\phantomsection\label{ux434ux43eux43aux43bux430ux434ux447ux438ux43a}
\begin{columns}[c]
\begin{column}{0.65\linewidth}
\begin{itemize}[<+->]
\tightlist
\item
  \textbf{Карпухин Клим}
\item
  Российский университет дружбы народов
\item
  Email:
  \href{mailto:1032255580@rudn.ru}{\nolinkurl{1032255580@rudn.ru}}
\item
  Роли: студент (лабораторная работа по ОС/виртуализации)
\end{itemize}
\end{column}

\begin{column}{0.35\linewidth}

\includegraphics[width=0.9\linewidth,height=\textheight,keepaspectratio]{./image/dokladchik.png}
\end{column}
\end{columns}
\end{frame}

\section{Вводная
часть}\label{ux432ux432ux43eux434ux43dux430ux44f-ux447ux430ux441ux442ux44c}

\begin{frame}{Актуальность}
\protect\phantomsection\label{ux430ux43aux442ux443ux430ux43bux44cux43dux43eux441ux442ux44c}
\begin{itemize}[<+->]
\tightlist
\item
  Современные специалисты ежедневно работают с множеством учётных
  записей и конфигурационных файлов. Безопасное хранение паролей и
  управление настройками рабочей среды --- критически важные задачи.
\item
  Инструменты pass (менеджер паролей) и chezmoi (управление dotfiles)
  позволяют решить эти задачи с использованием проверенных
  криптографических методов (GPG) и систем контроля версий (Git).
\item
  Синхронизация паролей и конфигураций между несколькими машинами
  становится простой и надёжной, что особенно актуально для
  разработчиков и администраторов.
\end{itemize}
\end{frame}

\begin{frame}[fragile]{Объект и предмет исследования}
\protect\phantomsection\label{ux43eux431ux44aux435ux43aux442-ux438-ux43fux440ux435ux434ux43cux435ux442-ux438ux441ux441ux43bux435ux434ux43eux432ux430ux43dux438ux44f}
\begin{itemize}[<+->]
\tightlist
\item
  \textbf{Объект:} процесс безопасного хранения учётных данных и
  управления конфигурационными файлами в Unix-подобных операционных
  системах.
\item
  \textbf{Предмет:} менеджер паролей \texttt{pass}, его расширения
  (\texttt{gopass}, \texttt{browserpass}), система управления dotfiles
  \texttt{chezmoi}, а также технологии GPG и Git как основа для их
  работы.
\end{itemize}
\end{frame}

\section{Научная новизна и практическая
значимость}\label{ux43dux430ux443ux447ux43dux430ux44f-ux43dux43eux432ux438ux437ux43dux430-ux438-ux43fux440ux430ux43aux442ux438ux447ux435ux441ux43aux430ux44f-ux437ux43dux430ux447ux438ux43cux43eux441ux442ux44c}

\begin{frame}[fragile]{Научная новизна}
\protect\phantomsection\label{ux43dux430ux443ux447ux43dux430ux44f-ux43dux43eux432ux438ux437ux43dux430}
\begin{itemize}[<+->]
\tightlist
\item
  Интеграция GPG-шифрования с файловой системой для организации
  иерархического хранилища паролей.
\item
  Использование Git для версионирования и синхронизации зашифрованных
  данных без раскрытия их содержимого.
\item
  Применение шаблонов (Go templates) в \texttt{chezmoi} для адаптации
  конфигураций под различные машины в рамках единого репозитория
  dotfiles.
\end{itemize}
\end{frame}

\begin{frame}{Практическая значимость работы}
\protect\phantomsection\label{ux43fux440ux430ux43aux442ux438ux447ux435ux441ux43aux430ux44f-ux437ux43dux430ux447ux438ux43cux43eux441ux442ux44c-ux440ux430ux431ux43eux442ux44b}
\begin{itemize}[<+->]
\tightlist
\item
  Полученные навыки позволяют организовать единое безопасное хранилище
  паролей, доступное с любого устройства, и автоматически разворачивать
  рабочее окружение на новых машинах.
\item
  Снижение риска потери или компрометации паролей за счёт шифрования и
  резервного копирования через Git.
\item
  Ускорение настройки новых рабочих мест и упрощение поддержки
  существующих конфигураций.
\end{itemize}
\end{frame}

\section{Цель, гипотеза и
задачи}\label{ux446ux435ux43bux44c-ux433ux438ux43fux43eux442ux435ux437ux430-ux438-ux437ux430ux434ux430ux447ux438}

\begin{frame}[fragile]{Цель}
\protect\phantomsection\label{ux446ux435ux43bux44c}
Освоить работу с менеджером паролей \texttt{pass} и системой управления
файлами конфигурации \texttt{chezmoi}. Получить навыки безопасного
хранения паролей, синхронизации их через Git, а также управления
конфигурационными файлами на нескольких машинах.
\end{frame}

\begin{frame}[fragile]{Гипотеза}
\protect\phantomsection\label{ux433ux438ux43fux43eux442ux435ux437ux430}
Применение \texttt{pass} и \texttt{chezmoi} позволит значительно
повысить безопасность и воспроизводимость рабочей среды, упростив её
перенос между устройствами.
\end{frame}

\begin{frame}[fragile]{Задачи}
\protect\phantomsection\label{ux437ux430ux434ux430ux447ux438}
\begin{itemize}[<+->]
\tightlist
\item
  Установить и настроить менеджер паролей \texttt{pass} и его расширение
  \texttt{gopass}.
\item
  Сгенерировать GPG-ключ и инициализировать хранилище паролей.
\item
  Подключить удалённый Git-репозиторий для синхронизации хранилища.
\item
  Установить браузерное расширение \texttt{browserpass} для удобного
  использования паролей.
\item
  Установить \texttt{chezmoi} и создать собственный репозиторий dotfiles
  на основе шаблона.
\item
  Протестировать управление конфигурацией на локальной машине и
  развернуть её на другой системе.
\end{itemize}
\end{frame}

\section{Материалы и
методы}\label{ux43cux430ux442ux435ux440ux438ux430ux43bux44b-ux438-ux43cux435ux442ux43eux434ux44b}

\begin{frame}{Материалы и методы}
\protect\phantomsection\label{ux43cux430ux442ux435ux440ux438ux430ux43bux44b-ux438-ux43cux435ux442ux43eux434ux44b-1}
\begin{itemize}[<+->]
\tightlist
\item
  \textbf{Инструменты:}

  \begin{itemize}[<+->]
  \tightlist
  \item
    Менеджер паролей \textbf{pass} (стандартный менеджер паролей Unix).
  \item
    Расширенная версия \textbf{gopass} (с дополнительными
    возможностями).
  \item
    Криптографическая система \textbf{GPG} для шифрования паролей.
  \item
    Система контроля версий \textbf{Git} и платформа \textbf{GitHub} для
    синхронизации.
  \item
    Браузерное расширение \textbf{browserpass} для интеграции с
    браузером.
  \item
    Инструмент управления dotfiles \textbf{chezmoi}.
  \item
    Редактор и утилиты командной строки (установка шрифтов, бинарных
    файлов).
  \end{itemize}
\end{itemize}
\end{frame}

\section{Ход работы --- установка и настройка
pass}\label{ux445ux43eux434-ux440ux430ux431ux43eux442ux44b-ux443ux441ux442ux430ux43dux43eux432ux43aux430-ux438-ux43dux430ux441ux442ux440ux43eux439ux43aux430-pass}

\begin{frame}[fragile]{Этап 1: Установка pass и gopass}
\protect\phantomsection\label{ux44dux442ux430ux43f-1-ux443ux441ux442ux430ux43dux43eux432ux43aux430-pass-ux438-gopass}
\begin{itemize}[<+->]
\tightlist
\item
  Установка менеджера паролей \texttt{pass} из репозитория.
\end{itemize}

\begin{figure}

\centering{

\includegraphics[width=0.3\linewidth,height=\textheight,keepaspectratio]{image/1.png}

}

\caption{\label{fig-001}Установка pass}

\end{figure}%

\begin{itemize}[<+->]
\tightlist
\item
  Установка расширенной версии \texttt{gopass}.
\end{itemize}

\begin{figure}

\centering{

\includegraphics[width=0.3\linewidth,height=\textheight,keepaspectratio]{image/2.png}

}

\caption{\label{fig-002}Установка gopass}

\end{figure}%
\end{frame}

\begin{frame}{Этап 2: Настройка GPG и инициализация хранилища}
\protect\phantomsection\label{ux44dux442ux430ux43f-2-ux43dux430ux441ux442ux440ux43eux439ux43aux430-gpg-ux438-ux438ux43dux438ux446ux438ux430ux43bux438ux437ux430ux446ux438ux44f-ux445ux440ux430ux43dux438ux43bux438ux449ux430}
\begin{itemize}[<+->]
\tightlist
\item
  Просмотр существующих GPG-ключей.
\end{itemize}

\begin{figure}

\centering{

\includegraphics[width=0.3\linewidth,height=\textheight,keepaspectratio]{image/3.png}

}

\caption{\label{fig-003}Просмотр ключей}

\end{figure}%

\begin{itemize}[<+->]
\tightlist
\item
  Инициализация хранилища паролей с использованием выбранного ключа.
\end{itemize}

\begin{figure}

\centering{

\includegraphics[width=0.3\linewidth,height=\textheight,keepaspectratio]{image/4.png}

}

\caption{\label{fig-004}Инициализация хранилища}

\end{figure}%
\end{frame}

\begin{frame}{Этап 3: Синхронизация через Git}
\protect\phantomsection\label{ux44dux442ux430ux43f-3-ux441ux438ux43dux445ux440ux43eux43dux438ux437ux430ux446ux438ux44f-ux447ux435ux440ux435ux437-git}
\begin{itemize}[<+->]
\tightlist
\item
  Создание Git-структуры внутри хранилища.
\end{itemize}

\begin{figure}

\centering{

\includegraphics[width=0.3\linewidth,height=\textheight,keepaspectratio]{image/5.png}

}

\caption{\label{fig-005}Создание структуры git}

\end{figure}%

\begin{itemize}[<+->]
\tightlist
\item
  Привязка удалённого репозитория на GitHub.
\end{itemize}

\begin{figure}

\centering{

\includegraphics[width=0.3\linewidth,height=\textheight,keepaspectratio]{image/6.png}

}

\caption{\label{fig-006}Задание адреса репозитория}

\end{figure}%
\end{frame}

\begin{frame}{Этап 4: Добавление пароля и настройка browserpass}
\protect\phantomsection\label{ux44dux442ux430ux43f-4-ux434ux43eux431ux430ux432ux43bux435ux43dux438ux435-ux43fux430ux440ux43eux43bux44f-ux438-ux43dux430ux441ux442ux440ux43eux439ux43aux430-browserpass}
\begin{itemize}[<+->]
\tightlist
\item
  Добавление нового пароля в хранилище.
\end{itemize}

\begin{figure}

\centering{

\includegraphics[width=0.3\linewidth,height=\textheight,keepaspectratio]{image/8.png}

}

\caption{\label{fig-008}Добавление нового пароля}

\end{figure}%

\begin{itemize}[<+->]
\tightlist
\item
  Настройка интерфейса для взаимодействия с браузером (browserpass).
\end{itemize}

\begin{figure}

\centering{

\includegraphics[width=0.3\linewidth,height=\textheight,keepaspectratio]{image/7.png}

}

\caption{\label{fig-007}Настройка интерфейса для браузера}

\end{figure}%
\end{frame}

\section{Ход работы --- установка и настройка
chezmoi}\label{ux445ux43eux434-ux440ux430ux431ux43eux442ux44b-ux443ux441ux442ux430ux43dux43eux432ux43aux430-ux438-ux43dux430ux441ux442ux440ux43eux439ux43aux430-chezmoi}

\begin{frame}{Этап 5: Установка chezmoi и создание репозитория dotfiles}
\protect\phantomsection\label{ux44dux442ux430ux43f-5-ux443ux441ux442ux430ux43dux43eux432ux43aux430-chezmoi-ux438-ux441ux43eux437ux434ux430ux43dux438ux435-ux440ux435ux43fux43eux437ux438ux442ux43eux440ux438ux44f-dotfiles}
\begin{itemize}[<+->]
\tightlist
\item
  Установка дополнительного ПО (chezmoi, шрифты, бинарные файлы).
\end{itemize}

\includegraphics[width=0.3\linewidth,height=\textheight,keepaspectratio,alt={Установка дополнительного ПО}]{image/9.png}
\includegraphics[width=0.3\linewidth,height=\textheight,keepaspectratio,alt={Установка шрифтов}]{image/10.png}

\begin{itemize}[<+->]
\tightlist
\item
  Создание собственного репозитория dotfiles с помощью утилит chezmoi.
\end{itemize}

\begin{figure}

\centering{

\includegraphics[width=0.3\linewidth,height=\textheight,keepaspectratio]{image/12.png}

}

\caption{\label{fig-012}Создание репозитория с помощью утилит}

\end{figure}%
\end{frame}

\begin{frame}{Этап 6: Подключение репозитория и работа на другой машине}
\protect\phantomsection\label{ux44dux442ux430ux43f-6-ux43fux43eux434ux43aux43bux44eux447ux435ux43dux438ux435-ux440ux435ux43fux43eux437ux438ux442ux43eux440ux438ux44f-ux438-ux440ux430ux431ux43eux442ux430-ux43dux430-ux434ux440ux443ux433ux43eux439-ux43cux430ux448ux438ux43dux435}
\begin{itemize}[<+->]
\tightlist
\item
  Подключение созданного репозитория к текущей системе.
\end{itemize}

\begin{figure}

\centering{

\includegraphics[width=0.3\linewidth,height=\textheight,keepaspectratio]{image/13.png}

}

\caption{\label{fig-013}Подключение репозитория к системе}

\end{figure}%

\begin{itemize}[<+->]
\tightlist
\item
  На второй машине инициализация chezmoi с тем же репозиторием dotfiles.
\end{itemize}

\begin{figure}

\centering{

\includegraphics[width=0.3\linewidth,height=\textheight,keepaspectratio]{image/14.png}

}

\caption{\label{fig-014}Инициализация chezmoi с репозиторием dotfiles}

\end{figure}%

\begin{itemize}[<+->]
\tightlist
\item
  Добавление параметров autoCommit и autoPush в конфигурацию для
  автоматической синхронизации.
\end{itemize}

\begin{figure}

\centering{

\includegraphics[width=0.3\linewidth,height=\textheight,keepaspectratio]{image/15.png}

}

\caption{\label{fig-015}Добавление autoCommit и autoPush в конфигурацию}

\end{figure}%
\end{frame}

\section{Результаты и
анализ}\label{ux440ux435ux437ux443ux43bux44cux442ux430ux442ux44b-ux438-ux430ux43dux430ux43bux438ux437}

\begin{frame}[fragile]{Анализ достигнутых результатов}
\protect\phantomsection\label{ux430ux43dux430ux43bux438ux437-ux434ux43eux441ux442ux438ux433ux43dux443ux442ux44bux445-ux440ux435ux437ux443ux43bux44cux442ux430ux442ux43eux432}
\begin{itemize}[<+->]
\tightlist
\item
  Создано зашифрованное хранилище паролей на базе \texttt{pass},
  инициализировано с GPG-ключом и связано с удалённым Git-репозиторием.
\item
  Пароли успешно добавляются, извлекаются и синхронизируются между
  устройствами.
\item
  Настроено браузерное расширение \texttt{browserpass}, позволяющее
  автоматически подставлять пароли на веб-сайтах.
\item
  С помощью \texttt{chezmoi} создан репозиторий dotfiles, который
  развёрнут на двух машинах; конфигурация автоматически применяется и
  отслеживается через Git.
\end{itemize}
\end{frame}

\begin{frame}[fragile]{Практическая значимость результатов}
\protect\phantomsection\label{ux43fux440ux430ux43aux442ux438ux447ux435ux441ux43aux430ux44f-ux437ux43dux430ux447ux438ux43cux43eux441ux442ux44c-ux440ux435ux437ux443ux43bux44cux442ux430ux442ux43eux432}
\begin{itemize}[<+->]
\tightlist
\item
  Полученная инфраструктура обеспечивает высокий уровень безопасности
  учётных данных и лёгкость восстановления окружения при смене
  оборудования.
\item
  Навыки работы с \texttt{chezmoi} позволяют эффективно поддерживать
  единообразие конфигураций в команде или на личных устройствах.
\item
  Достигнутая автоматизация (autoCommit/autoPush) минимизирует ручные
  операции при изменении конфигураций.
\end{itemize}
\end{frame}

\section{Выводы}\label{ux432ux44bux432ux43eux434ux44b}

\begin{frame}[fragile]{Общее заключение}
\protect\phantomsection\label{ux43eux431ux449ux435ux435-ux437ux430ux43aux43bux44eux447ux435ux43dux438ux435}
\begin{itemize}[<+->]
\tightlist
\item
  В ходе лабораторной работы освоены современные инструменты для
  безопасного хранения паролей (\texttt{pass}, \texttt{gopass}) и
  централизованного управления конфигурационными файлами
  (\texttt{chezmoi}).
\item
  На практике опробованы: генерация и использование GPG-ключей,
  инициализация и синхронизация хранилища паролей через Git, интеграция
  pass с браузером, создание собственного репозитория dotfiles,
  применение шаблонов и работа с chezmoi на нескольких машинах.
\item
  Цель работы достигнута, гипотеза подтверждена: использование данных
  инструментов делает рабочую среду безопаснее и воспроизводимее.
\end{itemize}
\end{frame}

\begin{frame}[fragile]{Выводы}
\protect\phantomsection\label{ux432ux44bux432ux43eux434ux44b-1}
\begin{enumerate}[<+->]
\tightlist
\item
  Менеджер паролей \texttt{pass} в сочетании с Git и GPG предоставляет
  простое, но надёжное решение для хранения и синхронизации паролей.
\item
  Расширение \texttt{browserpass} существенно упрощает повседневное
  использование паролей в браузере.
\item
  \texttt{Chezmoi} эффективно решает задачу управления dotfiles
  благодаря шаблонам и интеграции с Git.
\item
  Автоматическая фиксация и отправка изменений (autoCommit/autoPush)
  снижают риск рассинхронизации конфигураций.
\item
  Полученные компетенции могут быть непосредственно применены в учебной,
  научной и профессиональной деятельности для организации безопасной и
  переносимой рабочей среды.
\end{enumerate}
\end{frame}




\end{document}
